\documentclass[conference]{IEEEtran}
\IEEEoverridecommandlockouts
\usepackage[T1]{fontenc}
\usepackage[utf8]{inputenc}
\usepackage{cite}
\usepackage{amsmath,amssymb,amsfonts}
\usepackage{graphicx}
\usepackage{booktabs}
\usepackage{multirow}
\usepackage{textcomp}
\usepackage{xcolor}
\begin{document}

\title{LiteSemRAG: Lightweight Semantic-Aware Graph Retrieval for Robust RAG}

\author{}


\maketitle

\begin{abstract}

Graph-based Retrieval-Augmented Generation (RAG) improves structured evidence aggregation, but many graph-based RAG frameworks rely on large language models (LLMs) during indexing and querying, causing high token cost, latency, and deployment complexity. We propose \textbf{LiteSemRAG}, a lightweight semantic-aware graph retrieval framework that reduces repeated generative LLM dependence while preserving contextual semantic retrieval. LiteSemRAG constructs a heterogeneous semantic graph that separates lexical anchors, contextual semantic nodes, and text chunks to model polysemy. Semantic nodes are derived from contextual span representations and used as the central units for graph-based retrieval, while ambiguous lexical units can be associated with human-readable sense labels to support inspectable disambiguation. At query time, LiteSemRAG matches query semantics to the graph and ranks candidate chunks by integrating semantic relevance with local support among matched meanings. Experiments on three benchmarks show that LiteSemRAG achieves competitive or superior Recall@10 and MRR@10 compared with state-of-the-art LLM-based graph RAG systems, while substantially reducing indexing and query cost.

\end{abstract}

\begin{IEEEkeywords}
Retrieval-Augmented Generation, Graph-Based Retrieval, Heterogeneous Graph
\end{IEEEkeywords}

\section{Introduction}
Retrieval-Augmented Generation (RAG) \cite{lewis2020retrieval,borgeaud2022improving,izacard2021leveraging,zhao2026retrieval} has become a fundamental paradigm for improving factual grounding and knowledge integration in large language models (LLMs) \cite{brown2020language,devlin2019bert}.  By augmenting generation with retrieval over external corpora, RAG systems can improve factuality, adaptability, and domain generalization in knowledge-intensive tasks \cite{guu2020retrieval}. Therefore, the quality and efficiency of the retrieval component have become critical factors in RAG systems. Early RAG systems mainly rely on flat chunk-level retrieval using dense or sparse representations~\cite{karpukhin2020dense,khattab2020colbert}. However, RAG systems with flat chunk-level retrieval often struggle with document structure, semantic variation, and multi-hop evidence dependencies. To address these limitations, recent work has introduced structure-aware and graph-based retrieval frameworks \cite{sarthi2024raptor,gutierrez2024hipporag,edge2024local,guo2024lightrag}. They construct hierarchical or graph-based index structures to enable multi-level reasoning and structured evidence aggregation. These approaches demonstrate that structured indexing can substantially improve retrieval effectiveness, particularly in complex reasoning scenarios. However, many existing graph-based RAG frameworks heavily rely on LLMs during both indexing and querying. Although LLMs provide powerful capabilities for semantic extraction, summarization, and reasoning, utilization of LLMs also introduces substantial computational overhead, high token consumption, latency variability, and sensitivity to the specific underlying model. These factors can significantly limit the practicality of such systems in resource-constrained or large-scale scenarios. NoLLMRAG~\cite{yunollmrag} first attempts to eliminate LLM dependence in graph-based RAG. However, it only relies on surface-level lexical and static statistical structures, which limits contextual and semantic understanding.

In this paper, we propose \textbf{LiteSemRAG}, a lightweight and LLM-frugal semantic-aware graph retrieval framework. Unlike existing graph-based RAG frameworks that depend on LLMs for all semantic extraction and reasoning, LiteSemRAG demonstrates that strong semantic-aware indexing and retrieval can be achieved without making LLMs the recurring engine for graph construction or query-time reasoning. LiteSemRAG constructs a heterogeneous semantic graph directly by exploiting contextual embeddings. The heterogeneous semantic graph explicitly separates surface lexical representations (token nodes) from context-dependent semantic meanings (semantic nodes), enabling dynamic modeling of polysemy while avoiding LLM-centered graph construction. The index also retains multiword expressions as lexical evidence, helping to preserve phrase-level precision while modeling contextual meanings. To enhance robustness, LiteSemRAG applies an anchor-guided semantic assignment algorithm to words and phrases that appear with multiple contextual meanings. The assignment algorithm samples a small anchor set with farthest-first traversal and random coverage, propagates anchor labels over a mutual-kNN occurrence graph, uses cross-encoder prediction scores to veto unreliable neighbor overrides, and resolves remaining uncertain occurrences with cross-encoder fallback. LiteSemRAG only uses targeted LLM assistance to merge and filter candidate sense descriptions for semantic-label preparation, as it requires complicated sense consolidation and noise removal that conventional non-generative models cannot reliably perform through fixed-label classification or similarity scoring alone. At query stage, LiteSemRAG resolves query units through exact, partial, and embedding-based semantic-node matching, and ranks chunks by combining direct semantic evidence with bounded pairwise co-occurrence support. This design keeps retrieval a deterministic encoder-driven process while preserving both contextual semantic disambiguation and lexical precision.
We evaluate LiteSemRAG on three widely used benchmarks: HotpotQA, SciFact, and FEVER. Experimental results show that LiteSemRAG achieves competitive or superior Recall@10 and MRR@10 compared with two state-of-the-art LLM-based graph RAG systems. This design makes LiteSemRAG particularly suitable for large-scale retrieval systems where LLM-based indexing introduces prohibitive computational and monetary costs.
The contributions of this paper are summarized as follows:
\begin{itemize}
    \item We propose a lightweight and LLM-frugal semantic-aware graph retrieval framework that explicitly models polysemy through contextual embeddings.
    \item We introduce an anchor-guided semantic assignment algorithm, combining farthest-first anchor sampling, mutual-kNN label propagation, cross-encoder margin veto, and cross-encoder fallback to effectively assign semantics to samples.
    \item We design a group-aware semantic co-occurrence retrieval strategy that combines semantic-node matching, direct semantic evidence, and bounded pairwise chunk support.
    \item We demonstrate that LiteSemRAG achieves competitive or superior retrieval performance compared to LLM-based graph RAG systems while significantly improving efficiency through targeted indexing-time LLM use and encoder-driven retrieval.
\end{itemize}


\section{Related Work}

\subsection{Retrieval-Augmented Generation}

Retrieval-augmented generation (RAG) \cite{lewis2020retrieval} combines a parametric generator with a non-parametric retriever on an external corpus. RAG has become a core component for improving performance in knowledge-intensive natural language processing tasks. However, many RAG frameworks still rely on flat chunk retrieval, which can struggle to capture document structure, semantic variation, and multi-hop evidence dependencies \cite{karpukhin2020dense}. Some works are proposed to improve the retrieval component through dense or sparse neural representations. Karpukhin et al. \cite{karpukhin2020dense} proposed a dual-encoder dense retriever that maps questions and passages into a shared embedding space for nearest-neighbor search. ColBERT \cite{khattab2020colbert} is proposed by exploiting late interaction to preserve token-level matching signals while remaining efficient for large-scale retrieval. These methods substantially improve retrieval quality, but they generally retrieve over flat format collections and do not explicitly induce graph-structured semantic units.

% \subsection{Neural Retrieval Models}
% A major line of work improves the retrieval component itself through dense or sparse neural representations. Karpukhin et al. proposed a dual-encoder dense retriever that maps questions and passages into a shared embedding space for nearest-neighbor search \cite{karpukhin2020dense}. \emph{ColBERT: Efficient and Effective Passage Search via Contextualized Late Interaction over BERT} proposed late interaction to preserve token-level matching signals while remaining efficient enough for large-scale retrieval. :contentReference[oaicite:3]{index=3} In contrast, \emph{SPLADE: Sparse Lexical and Expansion Model for First Stage Ranking} learns sparse lexical representations that retain the efficiency and exact-matching benefits of inverted indexing. :contentReference[oaicite:4]{index=4} These methods substantially improve retrieval quality, but they generally retrieve over flat passage collections and do not explicitly induce graph-structured semantic units or separate surface lexical forms from context-dependent semantic senses. :contentReference[oaicite:5]{index=5}

\subsection{Structure-Aware and Graph-Based RAG}
Recent work has explored structured retrieval to better support long-context reasoning and multi-hop evidence aggregation. RAPTOR \cite{sarthi2024raptor} builds a hierarchical tree of clustered and summarized text units, enabling retrieval at multiple abstraction levels. GraphRAG \cite{edge2024local} constructs an entity-centric graph and uses LLM-based extraction and community summaries to support global and local query answering. HippoRAG \cite{gutierrez2024hipporag} combines knowledge graphs with personalized PageRank to improve multi-hop retrieval efficiency. LightRAG \cite{guo2024lightrag} further simplifies graph-based RAG with dual-level retrieval over low-level and high-level knowledge structures. These methods show that structured indexing can improve retrieval quality beyond flat chunk search, especially for complex reasoning tasks. However, the dominant graph-based RAG frameworks still depend on LLMs during indexing and querying, which increases latency, cost, and reproducibility sensitivity to the underlying model. NoLLMRAG \cite{yunollmrag} is the first graph-based RAG framework that eliminates reliance on LLMs during indexing and querying. However, due to absence of explicit semantic modeling, it only relies on surface-level lexical and
static statistical structures, which can significantly limit its effectiveness in complex or multi-hop scenarios.

\section{Approach}
LiteSemRAG is a lightweight semantic retrieval framework that constructs a heterogeneous semantic graph to preserve contextual semantic meaning while minimizing dependence on generative LLM calls. The framework consists of two main components: \textbf{Heterogeneous Semantic Graph Index}, where semantic-aware structures are derived primarily from contextual token-level embeddings, and \textbf{Semantic Co-occurrence Retrieval}, where query spans are matched to semantic nodes and chunks are ranked by a unified evidence score. Figure~\ref{query} illustrates the query workflow. LiteSemRAG first analyzes the query with the same phrase-aware representation used during indexing, resolves each query unit to exact, partial, or embedding-based semantic-node matches, and then ranks candidate chunks by combining direct semantic evidence with bounded pairwise co-occurrence support. This design gives every matched semantic node a direct retrieval path while using co-occurrence evidence as a controlled score boosting.
We define a multi-layer heterogeneous semantic graph as $G = (V, E)$, where the node set is defined as $V = D \cup C \cup S \cup T$. The node types are defined as:
\begin{itemize}
    \item \textbf{Document nodes} ($D$): each represents a document in the corpus.
    \item \textbf{Chunk nodes} ($C$): each represents a text chunk derived from a document.
    \item \textbf{Semantic nodes} ($S$): each represents a specific semantic meaning of a token node.
    \item \textbf{Token nodes} ($T$): each represents a unique lexical anchor. A token node may correspond to a single token, a lexicalized entity or a phrase.
\end{itemize}
The edge set $E$ consists of bidirectional edges of three types:
\begin{itemize}
    \item Document--Chunk edges: indicating that a chunk belongs to a document.
    \item Chunk--Semantic edges: indicating that a semantic meaning appears in a chunk.
    \item Semantic--Token edges: linking each semantic node to its corresponding lexical anchor.
\end{itemize}

\subsection{Heterogeneous Semantic Graph Index}
Given a document collection, LiteSemRAG first splits each document into chunks. From each chunk, important textual elements (entities, phrases, and tokens) are extracted by NLP tools and treated as candidate spans. Before embedding and graph insertion, phrase candidates pass through a phrase-aware routing layer that classifies them as lexicalized or compositional. Lexicalized spans remain complete lexical units. Compositional spans are represented by a coupled group containing the original phrase surface, its semantic head, and its modifiers. For each routed occurrence, LiteSemRAG keeps both the lexical anchor used for semantic matching and the original evidence span, and then computes a context-dependent span embedding using a token-level text encoder. This embedding encoder is efficient because one forward pass over a chunk can provide contextual representations for all spans in that chunk. Semantic nodes are dynamically constructed to represent distinct contextual meanings based on these contextual span embeddings. By contrast, semantic judgment over candidate meanings is handled by a separate expensive but more precise text encoder, which must be run for each individual span-context and candidate-meaning comparison. The final index graph is organized as a multi-layer heterogeneous semantic graph connecting documents, chunks, semantic nodes, and token nodes, while occurrence-level information preserves phrase surfaces and modifier constraints.

\subsubsection{Phrase-Aware Lexical Routing.}
Multiword phrases are not uniformly suitable as independent semantic units. Some phrases behave as lexicalized expressions whose meaning is tied to the complete phrase, while others are compositional expressions whose head carries the main semantic type and whose modifiers provide additional constraints. LiteSemRAG therefore separates these two cases during indexing. Lexicalized expressions and named entities are preserved as complete phrase anchors, so semantic discovery operates over the full phrase. This is appropriate for expressions such as \emph{machine learning}, \emph{New York}, and \emph{White House}, where decomposing the phrase would lose the intended unit of meaning.

For compositional phrases, LiteSemRAG creates a phrase--head relation rather than treating the full phrase as an isolated unit. For example, \emph{American director} is represented by the complete surface phrase \emph{American director}, the semantic head \emph{director}, and the modifier \emph{American}. The full phrase preserves exact surface evidence, while the head allows semantically related occurrences such as \emph{French director} and \emph{American director} to share an aggregation route. The modifier is retained as occurrence-level evidence and serves as a query-time constraint. This design prevents the graph from relying exclusively on sparse modifier-head combinations, but it also avoids discarding the original phrase.

Formally, each occurrence can be viewed as containing an indexed lexical anchor, an evidence span, and optional compositional-group metadata $o = (c,\; a,\; z,\; M,\; r,\; g)$, where $c$ is the chunk, $a$ is the lexical anchor used for semantic discovery or matching, $z$ is the original evidence span, $M$ is the set of modifiers attached to the occurrence, $r$ is the role of the occurrence within a compositional group, and $g$ identifies the group. For lexicalized expressions, $a$ and $z$ are the same phrase and no compositional group is required. For compositional phrases, the full phrase and the head form related anchors under the same group, while modifiers are retained as constraints on the occurrence. This representation allows semantic-node construction to operate on compact lexical anchors while preserving phrase-level evidence for interpretation and modifier-aware retrieval.

\subsubsection{Dynamic semantic node creation}

A key component of LiteSemRAG is the dynamic construction of semantic nodes from contextual span embeddings. Instead of assigning a fixed representation to each surface textual element, LiteSemRAG creates distinct semantic nodes that capture context-dependent meanings of the same indexed lexical anchor. This enables explicit modeling of polysemy without using LLM-based parsing. For each routed lexical anchor $t$, LiteSemRAG collects all of its contextual occurrences. Each occurrence contains the chunk in which the anchor appears, the corresponding evidence span, modifier information when applicable, and the local context in which the span appears. We denote the occurrence set as $O_t = \{o_1,o_2,\dots,o_n\}$. For each occurrence $o_i$, LiteSemRAG recomputes a contextual span embedding $e_i$ from its local context, so that $e_i$ represents the semantic position of indexed span $t$ in that specific context. Let
$E_t = \{ e_1, e_2, \dots, e_n \}$ denote the set of normalized contextual span embeddings for span $t$. To determine whether a span may exhibit multiple semantic meanings, we measure the concentration of its embedding distribution by computing the dispersion score (S-mean), defined as
\[
S\text{-mean}(t) = \frac{1}{n} \sum_{i=1}^{n} \cos(e_i, \bar{e}_t),
\]
where $\bar{e}_t = \frac{1}{n} \sum_{i=1}^{n} e_i$. A high S-mean value indicates that embeddings are tightly concentrated around a single semantic direction, while a low S-mean suggests multiple contextual semantic regions. To avoid unnecessary splitting of highly frequent but uninformative spans, we jointly consider the inverse document frequency (IDF). LiteSemRAG triggers semantic discovery for a span $t$ only if
$
IDF(t) > \tau_{idf} \quad \text{and} \quad S\text{-mean}(t) < \tau_{disp}
$,
where $\tau_{\text{idf}}$ and $\tau_{\text{disp}}$ are predefined thresholds. If this condition is not satisfied, the span is treated as single-semantic and one semantic node is created by averaging all contextual span embeddings.

\subsubsection{Anchor-Guided Semantic Assignment.}
When the S-mean criterion indicates that a span may be polysemous, LiteSemRAG performs anchor-guided semantic assignment rather than committing to geometry-only clustering or exhaustive semantic judgment for every occurrence. The design goal is to use a small number of explicitly judged anchor occurrences to steer label propagation, while using the cross-encoder as a conservative verifier when graph propagation conflicts with a confident semantic judgment. The first step is to select a small anchor set $A_t$ from $O_t$. A fixed fraction of anchors is selected by \emph{farthest-first traversal} (FFT), and the remaining anchors are sampled randomly from the non-selected occurrences. Starting from an initial occurrence, FFT repeatedly selects the occurrence whose embedding has the largest distance to its nearest already selected embedding:
\[
o^* = \arg\max_{o_i \in O_t \setminus A_t} \min_{o_j \in A_t} d(e_i,e_j),
\]
where $d(e_i,e_j)=1-\cos(e_i,e_j)$. FFT encourages anchors to cover boundary and minority regions in the embedding space, while random supplementation reduces deterministic sampling bias and gives dense regions a chance to contribute anchors.

\subsubsection{Anchor Labeling and Candidate Meanings.}
For the routed lexical anchor $t$, LiteSemRAG retrieves candidate meanings $M_t = \{m_1,m_2,\dots,m_r\}$ from Wikidata according to the surface form of the span. The raw candidate list is merged and filtered with targeted LLM assistance to remove redundant or noisy descriptions before occurrence-level assignment. This preparation step is intentionally separated from occurrence-level scoring: Wikidata candidates often contain overlapping entities, alias-driven duplicates, overly broad definitions, and descriptions that require semantic rewriting or consolidation. Traditional non-LLM models can rank or compare given descriptions, but they are not well suited to producing a cleaned candidate-sense inventory from noisy open-world descriptions. Each retained candidate meaning is represented by its textual description and alias information. For every anchor occurrence $o_i \in A_t$, a semantic judge compares the local occurrence context with each retained candidate meaning and assigns the best semantic label $y_i = \arg\max_{m \in M_t} f_{P}(o_i,t,m)$, where $f_{P}(\cdot)$ denotes the score from the selected semantic judge. In our implementation, this occurrence-level judge is the cross-encoder used for semantic-description scoring. Because candidate meanings are retrieved from an external knowledge source and consolidated into explicit textual descriptions, the labels assigned in this multi-sense branch are human-readable sense descriptions rather than opaque cluster IDs. These descriptions are used to inspect the disambiguated senses of polysemous anchors, while semantic nodes outside this branch remain represented by their lexical anchors and contextual embeddings.

\subsubsection{Consensus-Based Early Exit.}
After anchor labeling, LiteSemRAG first checks whether the anchors already form a strong semantic consensus. Let $n_y$ denote the number of anchors assigned to candidate meaning $y$. If $\max_y n_y / |A_t| \geq \tau_{\text{cons}}$, where $\tau_{\text{cons}}$ is a consensus threshold, LiteSemRAG treats the span as effectively single-semantic in the current corpus and assigns the majority anchor label to all occurrences of $t$. This early-exit rule avoids unnecessary semantic splitting when the anchor evidence indicates that nearly all observed contexts share the same meaning.

\subsubsection{Mutual-kNN Label Propagation.}
If no strong consensus is observed, the span is treated as potentially multi-semantic. LiteSemRAG builds a mutual-kNN graph over the normalized occurrence embeddings. Let $\mathcal{N}_k(i)$ denote the $k$ nearest neighbors of occurrence $o_i$ under cosine distance. An undirected edge is created only when the neighbor relation is reciprocal:
\[
i \sim j \quad \Longleftrightarrow \quad
j \in \mathcal{N}_k(i) \;\; \text{and} \;\; i \in \mathcal{N}_k(j).
\]
This mutual-neighbor constraint removes many asymmetric edges that arise from hubs or sparse boundary points, making propagation more conservative than plain kNN voting.

Propagation proceeds in rounds from the labeled anchors. For an unlabeled occurrence $o_u$, LiteSemRAG collects the labels of its currently labeled mutual-kNN neighbors. If the majority label is $y$ and its vote ratio is
\[
\rho_u(y)=
\frac{|\{o_j \in \mathcal{N}^{\text{mut}}_k(u): y_j=y\}|}
{|\{o_j \in \mathcal{N}^{\text{mut}}_k(u): y_j \text{ is assigned}\}|},
\]
then $y$ can be propagated when $\rho_u(y)$ exceeds the required vote threshold. This iterative procedure expands anchor labels through locally coherent embedding neighborhoods without invoking a semantic judge for every easy occurrence.

\subsubsection{Cross-Encoder Margin Veto.}
LiteSemRAG augments mutual-kNN voting with cross-encoder margin evidence. It computes cross-encoder scores between each occurrence and every retained candidate meaning. Let $y^{\text{CE}}_u$ be the cross-encoder top-1 label for occurrence $o_u$, and let
\[
\Delta_u = f_P(o_u,t,y^{\text{CE}}_u) -
\max_{y \neq y^{\text{CE}}_u} f_P(o_u,t,y)
\]
be the top-1 margin. If neighbor voting proposes a label that conflicts with $y^{\text{CE}}_u$, and the cross-encoder margin is high while the proposed label has a substantially lower score, the cross-encoder veto blocks the propagation decision. This prevents a dense local neighborhood from overwriting cases where the semantic judge is strongly confident. Conversely, when the cross-encoder margin is low, the method allows locally consistent neighbor evidence to override the cross-encoder top-1 label. The assignment rule therefore uses graph propagation mainly for ambiguous cases, while preserving confident semantic judgments when the candidate-level evidence is clear.

\subsubsection{Uncertain Cases and Semantic-Node Creation.}
After propagation stops, some occurrences may remain unlabeled because they have no reliable labeled mutual neighbors, because their vote ratio is too weak, or because the cross-encoder veto blocks their propagated label. These uncertain occurrences are resolved by cross-encoder fallback, i.e., $y_u = y^{\text{CE}}_u$. Thus, anchor propagation is allowed to correct low-confidence or locally inconsistent cross-encoder decisions, but the system still has a deterministic semantic fallback for cases where graph evidence is insufficient. After all occurrences are assigned semantic labels, LiteSemRAG creates one semantic node for each retained meaning. The anchor embedding of semantic node $s_y$ is computed by averaging all contextual embeddings assigned to $y$, $e_{s_y} = |O_y|^{-1} \sum_{o_i \in O_y} e_i$, where $O_y$ denotes the set of occurrences assigned to meaning $y$. Chunk nodes are then linked to the corresponding semantic nodes according to the occurrence-level assignments, and the original token node remains connected to all of its semantic nodes as a lexical anchor. Because each occurrence retains its evidence span, a semantic edge can still expose the original phrase surface and modifiers even when semantic discovery was performed on a head token. The final graph therefore keeps both compact semantic anchors and phrase-level evidence.

\subsubsection{Design Rationale}
This dynamic semantic node creation process allows LiteSemRAG to: 1) explicitly model polysemy without LLM-centered graph construction, 2) avoid committing semantic-node boundaries to embedding geometry alone, 3) use sparse anchor supervision to correct low-confidence cross-encoder confusions, and 4) maintain robustness by combining mutual-kNN locality with cross-encoder fallback. The resulting assignment strategy provides a compact speed-accuracy tradeoff. FFT and random sampling expose a small but diverse set of anchors; mutual-kNN propagation limits label spread to reciprocal local neighborhoods; the cross-encoder veto prevents aggressive neighbor overrides; and CE fallback gives every remaining occurrence a semantic assignment. In addition, retrieving candidate meanings from an external knowledge base provides explicit textual descriptions for the semantic labels used in polysemy resolution, so the resulting multi-sense assignments can be inspected and validated. LiteSemRAG therefore separates surface lexical representation from contextual semantic representation while using semantic nodes as retrieval units, with textual explanations reserved for the disambiguated senses produced by the polysemy-resolution process.

\subsection{Semantic Co-occurrence Retrieval}
At query stage, LiteSemRAG mirrors the phrase-aware index representation and ranks chunks with a unified semantic evidence score. The goal is to preserve lexical precision for exact phrase evidence, use semantic nodes to handle contextual meaning, and reward chunks where multiple query-relevant meanings provide local support for each other. Given a query, LiteSemRAG first constructs query units, resolves these units to semantic-node candidates, computes direct chunk evidence, and finally applies a bounded co-occurrence boost over semantic-node pairs that appear in the same candidate chunk.

\subsubsection{Composition-Aware Query Unit Construction}
For lexicalized query spans, the query unit remains the complete surface phrase. For compositional query spans, LiteSemRAG constructs a query group consisting of the full phrase, its semantic head, and its modifier units. For example, the query span \emph{American director} is represented by a full-phrase unit \emph{American director}, a head unit \emph{director}, and a modifier unit \emph{American}. The full phrase preserves exact lexical evidence when such a phrase is present in the index. The head unit supports aggregation over semantically related occurrences that share the same head. The modifier unit provides an auxiliary constraint for distinguishing surface realizations under the same head.

Let a compositional query group be denoted as $g_q = (p_q,\; h_q,\; M_q)$, where $p_q$ is the full phrase, $h_q$ is the semantic head, and $M_q$ is the set of modifier units. Each member of the group is assigned a role weight $\beta_r$, with $\beta_{\text{phrase}} > \beta_{\text{head}} > \beta_{\text{modifier}}$. Standalone query tokens whose character spans are already covered by a compositional phrase are not treated as independent query targets. This avoids double counting the same evidence while still allowing the full phrase, head, and modifiers to contribute through the shared group.

\subsubsection{Multi-Level Semantic Matching}
For each query unit $u_q$ with contextual embedding $e_q$, LiteSemRAG retrieves candidate \textbf{semantic nodes} rather than directly retrieving chunks. Matching is performed through three complementary routes:

\begin{itemize}
    \item \textbf{Exact match:} semantic nodes attached to the same token or phrase anchor as $u_q$.
    
    \item \textbf{Partial match:} semantic nodes attached to indexed phrase anchors that share the required phrase-level lexical evidence with $u_q$.
    
    \item \textbf{Similarity match:} semantic nodes whose embeddings are closest to $e_q$ when lexical matching does not provide sufficient candidates.
\end{itemize}

For exact matching, if the matched lexical anchor has multiple described semantic nodes, LiteSemRAG selects the semantic node whose textual description best matches the query context. If description-level disambiguation is unavailable or inconclusive, the selected semantic node is $s^*(t) = \arg\max_{s \in S(t)} \cos(e_q, e_s)$, where $S(t)$ denotes the semantic nodes linked to lexical anchor $t$, and $e_s$ is the semantic-node embedding. Partial matching follows the same semantic-node selection principle after retrieving compatible phrase anchors. Similarity matching is used as a fallback over the semantic-node embedding space and contributes a confidence proportional to its cosine similarity. These routes produce a set of matched semantic nodes $S_q$, where each semantic node is associated with an effective query weight $\eta_s$ determined by match confidence and, for compositional query units, by the role weight of the phrase, head, or modifier member.

\subsubsection{Direct Semantic Evidence}
The first component of chunk scoring measures how strongly a candidate chunk directly supports the matched query semantics. For a semantic node $s \in S_q$ and a chunk $c$, LiteSemRAG uses a semantic BM25-style contribution:
\[
\mathrm{BM25}(s,c) =
G(s)
\cdot
\frac{f(s,c) (k_1 + 1)}
{f(s,c) + k_1 \left(1 - b + b \frac{|c|}{\mathrm{avgcl}}\right)},
\]
where $G(s)$ is the chunk-level discriminative weight of semantic node $s$, $f(s,c)$ is the frequency of $s$ in chunk $c$, $|c|$ is the chunk length, and $\mathrm{avgcl}$ is the average chunk length. The direct semantic evidence score is then $\mathrm{Base}(c) = \sum_{s \in S_q(c)} \eta_s \cdot \mathrm{BM25}(s,c)$, where $S_q(c)$ is the set of matched semantic nodes appearing in chunk $c$. This score gives every matched semantic node a direct retrieval path, including nodes that do not co-occur with other query nodes.

For compositional query groups, direct evidence uses a group-aware aggregation rule. If a chunk contains the full-phrase semantic evidence for a group, the phrase contribution is used as the primary group contribution and the corresponding head or modifier members from the same group are not counted again. If the full phrase is absent, head and modifier units can jointly contribute according to their role weights. This rule reflects the asymmetric importance of phrase, head, and modifier evidence: exact phrase evidence is strongest, head evidence supports semantic aggregation, and modifier evidence refines the intended surface realization.

\subsubsection{Bounded Pairwise Co-occurrence Boosting}
The second component of chunk scoring rewards chunks where multiple matched semantic nodes support each other. For semantic nodes $i$ and $j$, LiteSemRAG defines a corpus-level co-occurrence strength:
\[
w_{ij} =
\frac{|C(i) \cap C(j)|}
{\sqrt{|C(i)| \cdot |C(j)|}},
\]
where $C(i)$ denotes the set of chunks containing semantic node $i$. This normalized overlap suppresses globally frequent semantics while emphasizing pairs that co-occur consistently across the corpus.

For a candidate chunk $c$, a semantic-node pair contributes a local boost only when both endpoints appear in $c$. The local evidence level $\ell_c(i,j)$ distinguishes stronger sentence-level support from weaker chunk-level support:
\[
\ell_c(i,j) =
\begin{cases}
\ell_{\text{sent}}, & \text{same sentence in } c,\\
\ell_{\text{chunk}}, & \text{same chunk but different sentences}.
\end{cases}
\]
Pairs formed only by different members of the same compositional query group are not used for this boost, because their phrase-level relation is already represented by the direct group-aware evidence. The pair contribution is therefore $\pi_c(i,j) = w_{ij} \cdot \ell_c(i,j)$. To prevent a long chunk or a dense semantic region from dominating the score through many weak pairs, LiteSemRAG averages only the top-$K$ pair contributions in each chunk:
\[
\mathrm{PairBoost}(c) =
\frac{1}{|\mathcal{P}_K(c)|}
\sum_{(i,j) \in \mathcal{P}_K(c)}
\pi_c(i,j),
\]
where $\mathcal{P}_K(c)$ contains the highest-scoring semantic-node pairs in chunk $c$ and $\mathrm{PairBoost}(c)=0$ when no valid pair is available.

\subsubsection{Final Chunk Ranking}
The final retrieval score combines direct evidence and bounded co-occurrence support as $\mathrm{Score}(c) = \mathrm{Base}(c) \cdot (1 + \lambda \cdot \mathrm{PairBoost}(c))$, where $\lambda$ controls the strength of the co-occurrence boost. This multiplicative form treats direct semantic evidence as the primary signal and uses co-occurrence as supporting evidence. A chunk that contains a strong individual semantic match can still rank through $\mathrm{Base}(c)$, while a chunk that jointly supports several query-relevant meanings receives an additional boost. Compared with a hard two-stage retrieval strategy, this unified ranking avoids separating pair-supported and single-node evidence into different retrieval pools; all candidate chunks are ranked by the same score, and co-occurrence improves the rank only when it is locally verified inside the candidate chunk.

\subsection{Extensibility of the Index Graph}
In the current framework, we restrict retrieval to direct semantic evidence and bounded co-occurrence-aware ranking to maintain computational efficiency and robustness. However, the underlying structure of LiteSemRAG's heterogeneous semantic index graph naturally supports more complex graph-based reasoning operations. The heterogeneous graph structure connects documents, chunks, semantic nodes, and token nodes in a bidirectional manner. This design enables flexible traversal across different elements and levels, allowing more advanced operations, including multi-hop semantic expansion, constrained subgraph extraction, and path-based relevance scoring. Therefore, LiteSemRAG should be viewed not only as a retrieval algorithm but also as a structured semantic indexing framework. The index graph provides a foundation for future extensions that incorporate deeper graph traversal strategies or reasoning mechanisms without requiring LLM-based graph construction.

\section{Experiments}
We evaluate LiteSemRAG against representative LLM-based graph RAG systems in terms of retrieval effectiveness and computational efficiency. Our experiments aim to answer the following questions:
\begin{itemize}
    \item Does LiteSemRAG achieve competitive retrieval performance compared to LLM-based graph RAG systems?
    \item How does LiteSemRAG compare in indexing and querying efficiency?
\end{itemize}
\subsection{Experimental Setup}
\subsubsection{Dataset}
We conduct experiments on three widely used retrieval benchmarks from the BEIR dataset \cite{thakur2021beir}: \textbf{HotpotQA} \cite{yang2018hotpotqa}: A multi-hop question answering dataset requiring reasoning across multiple documents; \textbf{SciFact} \cite{Wadden2020FactOF}: A scientific fact verification dataset focusing on evidence retrieval; \textbf{FEVER}: A large-scale fact verification dataset where claims must be supported or refuted using evidence sentences retrieved from Wikipedia.
For all datasets, documents are split into fixed-length chunks using the same preprocessing strategy across all compared methods to ensure fairness.
\subsubsection{Baseline and Setup}
We compare LiteSemRAG with two state-of-the-art LLM-based graph RAG systems: \textbf{GraphRAG} and \textbf{LightRAG}. Since they require LLMs for indexing and querying, we evaluate them under two LLM settings: \textbf{GPT-4o-mini} (API-based) and \textbf{LLaMA3-8B-Instruct} (locally deployed).
The two LLM-based RAG systems use identical LLM configurations within each setting. Both GraphRAG and LightRAG utilize \textbf{text-embedding-3-large} as the text embedding encoder. While LiteSemRAG uses \textbf{deberta-v3-large} as the token-level text encoder. The encoder difference arises because LiteSemRAG requires token-level contextual embeddings from a token-level text encoder, whereas the compared systems rely on sentence-level embeddings derived from sentence-level encoders. In the semantic discovery stage, LiteSemRAG further uses a more precise text encoder for semantic judgment over candidate meanings; in our experiments, we choose \textbf{cross-encoder/nli-deberta-v3-large} for anchor labeling, candidate-margin veto, and fallback assignment for unresolved occurrences during indexing. LiteSemRAG uses targeted LLM assistance to merge and filter Wikidata-derived candidate sense descriptions during semantic-label preparation, but the evaluated retrieval path does not depend on repeated generative LLM calls for graph construction or query-time reasoning. All experiments are conducted on a server equipped with two AMD EPYC 9354 CPUs and one AMD MI210 GPU.

\subsubsection{Evaluation Scope}
LiteSemRAG is designed as a retrieval framework rather than an end-to-end RAG system. Therefore, we evaluate only the quality of retrieved chunks. For fair comparison, we directly measure retrieval effectiveness using Recall@10 and MRR@10 based on ground-truth relevant chunks. We do not compare final answer quality obtained by feeding retrieved chunks into LLMs, since such evaluation would introduce additional variability from generation models and obscure differences in retrieval performance. This design isolates the retrieval component and ensures that performance differences reflect retrieval capability rather than downstream LLM reasoning or generation quality.
\subsubsection{Evaluation Metrics}
We evaluate retrieval quality by using \emph{Recall@10} and \emph{MRR@10}. For each query $q$ in a query dataset $Q$, let $R_q^{(10)}$ denote the ranked top-10 retrieved chunks, and let $G_q$ denote the set of chunks of gold (ground-truth) documents associated with $q$. In HotpotQA dataset, each query has exactly two gold documents. In SciFact and Fever datasets, each query has between one and five gold documents. We compute Recall@10 using a micro-averaged formulation:
\[
\text{Recall@10} =
\frac{
\sum_{q \in Q} | R_q^{(10)} \cap G_q |
}{
\sum_{q \in Q} | G_q |
}.
\]
This measures the proportion of gold documents that are successfully retrieved within the top 10 results across the entire dataset.

For each query $q$ in a query dataset $Q$, let $\text{rank}_q$ denote the rank position of the first retrieved document that belongs to $G_q$ within the top 10 results. If no gold document appears in the top 10, the reciprocal rank is defined as 0; otherwise, $\text{RR@10}(q)=1/\text{rank}_q$. The Mean Reciprocal Rank at 10 is computed as $\text{MRR@10}=|Q|^{-1}\sum_{q \in Q}\text{RR@10}(q)$. Recall@10 measures how completely the required evidence documents are retrieved, while MRR@10 evaluates how early at least one relevant document appears in the ranked chunk list. In addition, we report four efficiency metrics: Average Indexing Time per document (AIT), Average Query Time (AQT), average online generative-LLM Index Token Consumption per document (ITC), and average online generative-LLM Query Token Consumption (QTC). These metrics quantify both computational and cost efficiency.

\subsection{Retrieval Performance}
\begin{table*}[t]
\centering
\small
\setlength{\tabcolsep}{5pt}
\caption{Retrieval results on three datasets}
\label{retrieval results}
\begin{tabular}{lcccccc}
\hline
\multicolumn{1}{c}{\multirow{2}{*}{Model}} & \multicolumn{2}{c}{HotpotQA}     & \multicolumn{2}{c}{SciFact}      & \multicolumn{2}{c}{Fever}        \\ \cline{2-7}
\multicolumn{1}{c}{}                       & Recall@10       & MRR@10         & Recall@10       & MRR@10         & Recall@10       & MRR@10         \\ \hline
GraphRAG (GPT)                             & 55.6\%          & N/A            & 88.3\%          & N/A            & 44.9\%          & N/A            \\
GraphRAG (LLaMA)                           & ETL             & N/A            & ETL             & N/A            & ETL             & N/A            \\
LightRAG (GPT)                             & 72.1\%          & 0.429          & \textbf{95.4\%} & 0.419          & 84.5\%          & 0.497          \\
LightRAG (LLaMA)                           & 70.1\%          & 0.367          & 95.1\%          & 0.358          & 84.9\%          & 0.454          \\
LiteSemRAG (GPT)                           & \textbf{75.2\%} & 0.701          & 93.1\%          & \textbf{0.687} & \textbf{89.0\%} & \textbf{0.819} \\
LiteSemRAG (LLaMA)                         & 75.1\%          & 0.681          & 93.0\%          & 0.674          & 88.7\%          & 0.805           \\
NoLLMRAG                                   & 74.9\%          & \textbf{0.754} & 66.1\%          & 0.562          & 88.7\%          & 0.789
\end{tabular}
\end{table*}
Table~\ref{retrieval results} reports Recall@10 and MRR@10 on three datasets. Since GraphRAG does not provide explicit ranking of retrieved chunks, MRR@10 cannot be computed for GraphRAG. Therefore, we report only Recall@10 for GraphRAG. Additionally, when using LLaMA3-8B-Instruct as the underlying LLM, GraphRAG consistently encountered \textbf{Exceed maximum Token Limitation} (ETL) errors across all datasets, even though the chunk size was very small and contained only a few sentences (the maximum context length of LLaMA3-8B-Instruct is 8192 tokens). This error occurs not only during the indexing stage but also when querying databases that were already indexed using GPT. It prevented successful evaluation under this setting. Consequently, we mark these entries as ETL and exclude them from the comparison. The ETL results reflect the heavy reliance of GraphRAG on LLM-based operations during indexing and querying.

The main retrieval comparison focuses on Recall@10 and MRR@10. LiteSemRAG (GPT) achieves the strongest overall retrieval behavior among the semantic graph RAG methods: it obtains the best Recall@10 on HotpotQA and FEVER, remains close to the best Recall@10 on SciFact, and consistently provides much higher MRR@10 than LightRAG. This indicates that LiteSemRAG not only retrieves relevant evidence within the top results, but also ranks relevant chunks earlier. NoLLMRAG obtains strong scores on HotpotQA and FEVER, but its SciFact Recall@10 drops substantially. This pattern suggests that surface-level lexical and statistical matching can be effective when benchmark evidence is lexically aligned, but it is less stable when contextual semantic matching is required. The currently missing LiteSemRAG (LLaMA) results on SciFact and FEVER do not affect this conclusion: the complete LiteSemRAG (GPT) setting covers all three datasets, and the available HotpotQA LiteSemRAG (LLaMA) result remains close to the GPT setting and above the LLM-based graph baselines in both Recall@10 and MRR@10.

\subsection{Efficiency Analysis}
We report indexing and query efficiency on the three datasets in Table~\ref{tab:Time_results} and Table~\ref{tab:token_results}. Table~\ref{tab:Time_results} compares methods that perform explicit semantic analysis during graph-based retrieval. We do not include NoLLMRAG in this timing table because NoLLMRAG does not contain a semantic analysis module; its embeddings mainly support keyword-like matching, leading to extremely low AIT and AQT values that are not directly comparable with semantic graph RAG systems. Therefore, the relevant efficiency question is whether a method that still performs semantic understanding can reduce latency by design. LiteSemRAG addresses this question by constructing semantic nodes with encoder-driven representations and avoiding repeated online generative-LLM calls, rather than by removing semantic modeling altogether.

Under this comparable semantic-retrieval setting, LiteSemRAG achieves the lowest AIT and AQT across all three datasets. The improvement is especially clear at query time: GraphRAG remains expensive because its query workflow depends on LLM-centered operations, while LightRAG is faster but still slower than LiteSemRAG. LiteSemRAG reduces this latency by making retrieval an encoder-driven graph matching and co-occurrence scoring process after indexing. The missing LiteSemRAG (LLaMA) timing entries indicate unfinished measurements rather than a failed run, and they do not change the main efficiency conclusion because LiteSemRAG's evaluated retrieval path is not driven by the online choice of a generative LLM. Table~\ref{tab:token_results} further shows that LiteSemRAG confines generative-LLM usage to a small amount of indexing-stage semantic-description processing and completely removes online generative-LLM calls from query-time retrieval.

\begin{table*}[t]
\centering
\small
\setlength{\tabcolsep}{6pt}
\caption{Indexing and query time efficiency comparison across datasets}
\label{tab:Time_results}

\begin{tabular}{lcccccc}
\hline
                  & \multicolumn{2}{c}{HotpotQA}    & \multicolumn{2}{c}{SciFact}     & \multicolumn{2}{c}{Fever}       \\ \cline{2-7}
Method            & AIT            & AQT            & AIT            & AQT            & AIT            & AQT            \\ \hline
GraphRAG (GPT)    & 3.87s          & 29.9s          & 4.33s          & 51.21s         & 3.52s          & 30.17s         \\
GraphRAG (LLaMA)  & ETL            & ETL            & ETL            & ETL            & ETL            & ETL            \\
LightRAG (GPT)    & 8.14s          & 0.54s          & 10.54s         & 1.49s          & 7.76s          & 1.59s          \\
LightRAG (LLaMA)  & 8.36s          & 1.77s          & 10.59s         & 1.28s          & 8.35s          & 1.27s          \\
LiteSemRAG(GPT)   & \textbf{0.61s} & \textbf{0.09s} & \textbf{1.54s} & \textbf{0.10s} & \textbf{0.68s} & \textbf{0.06s} \\
LiteSemRAG(LLaMA) & 0.89s          & \textbf{0.09s} & 1.71s          & \textbf{0.10s} & 0.93s          & \textbf{0.06s} \\ \hline
\end{tabular}
\end{table*}


\begin{table*}[t]
\centering
\small
\setlength{\tabcolsep}{6pt}
\caption{Online generative-LLM token consumption comparison across datasets}
\label{tab:token_results}
\begin{tabular}{lcccccc}
\hline
                 & \multicolumn{2}{c}{HotpotQA} & \multicolumn{2}{c}{SciFact} & \multicolumn{2}{c}{Fever} \\ \cline{2-7} 
Method           & ITC           & QTC          & ITC          & QTC          & ITC         & QTC         \\ \hline
GraphRAG (GPT)   & 8.02K         & 10.85K       & 9.92K        & 11.24K       & 7.11K       & 4.06K       \\
GraphRAG (LLaMA) & ETL           & ETL          & ETL          & ETL          & ETL         & ETL         \\
LightRAG (GPT)   & 7.44K         & 0.43K        & 9.17K        & 0.59K        & 8.82K       & 0.44K       \\
LightRAG (LLaMA) & 8.17K         & 0.47K        & 8.85K        & 0.56K        & 8.23K       & 0.45K       \\
LiteSemRAG (GPT)   & \textbf{0.16K}    & \textbf{0}   & 0.46K   & \textbf{0}   & \textbf{0.18K}  & \textbf{0}  \\
LiteSemRAG (LLaMA) & 0.18K    & \textbf{0}   & \textbf{0.39K}   & \textbf{0}   & 0.21K  & \textbf{0}  \\ \hline
\end{tabular}

\end{table*}
Table~\ref{tab:token_results} reports the average online generative-LLM token consumption of the indexing and query stages. The main pattern is that LLM-based graph RAG baselines require substantial generative-LLM tokens during indexing, and GraphRAG also keeps a high token cost during querying because its retrieval workflow still invokes LLM-centered operations. LightRAG reduces query-time token usage compared with GraphRAG, but it does not remove generative-LLM dependence from the retrieval pipeline.

LiteSemRAG shows a different token-consumption profile. Its indexing-stage token usage is small because generative-LLM calls are limited to targeted semantic-description processing rather than broad graph extraction or reasoning. More importantly, LiteSemRAG's QTC is 0 on all three datasets for both GPT and LLaMA settings, indicating that once the semantic graph has been indexed, retrieval is fully encoder-driven and does not require online generative-LLM calls. Therefore, the token results support the main design claim: LiteSemRAG preserves semantic-aware indexing while replacing repeated generative-LLM dependence with a small, bounded indexing-stage cost and zero query-time generative-LLM cost.


\subsection{Ablation Study}
To evaluate the contribution of the semantic co-occurrence retrieval strategy, we conduct an ablation study using a broad search and reranking retrieval variant. In this setting, we first retrieve all semantic nodes matched by the query. All chunks connected to these semantic nodes are then collected, and a \textbf{jina-reranker-v3} reranker is applied to rank the candidate chunks. If the number of candidate chunks becomes extremely large, we randomly sample 100 chunks as reranking candidates. The top-10 chunks are selected as the final results. This ablation removes the bounded semantic co-occurrence scoring mechanism while preserving the same semantic index graph and matching process. In particular, both methods use the identical index graph constructed during the indexing stage and rely on the same token-level text encoder to generate contextual embeddings. The only difference lies in the retrieval strategy applied at query time. Table~\ref{tab:ablation_results} compares the standard LiteSemRAG and the broad search and reranking variant in terms of Recall@10 (R@10), MRR@10 (M@10), and Average Query Time (AQT). Since both methods operate on the same index graph, the comparison isolates the effect of the retrieval strategy itself. The results show that the proposed semantic co-occurrence retrieval significantly outperforms the broad retrieval baseline, demonstrating that the performance gains of LiteSemRAG primarily arise from group-aware semantic matching and locally verified co-occurrence scoring rather than the use of the token-level text encoder alone.

\begin{table*}[t]
\centering
\small
\setlength{\tabcolsep}{5pt}
\caption{Ablation study on semantic co-occurrence retrieval}
\label{tab:ablation_results}
\begin{tabular}{cccccccccc}
\hline
\multirow{2}{*}{Method}  & \multicolumn{3}{c}{HotpotQA} & \multicolumn{3}{c}{SciFact} & \multicolumn{3}{c}{Fever} \\ \cline{2-10} 
                         & R@10     & M@10    & AQT     & R@10     & M@10    & AQT    & R@10    & M@10   & AQT    \\ \hline
Broad search and reranking & 32.5\%   & 0.197   & 2.85s   & 14.9\%   & 0.056   & 2.97s  & 46.7\%  & 0.251  & 3.24s  \\
Semantic co-occurrence  & 75.2\%   & 0.701   & 0.09s   & 93.1\%   & 0.668   & 0.10s  & 89.0\%  & 0.819  & 0.06s  \\ \hline
\end{tabular}
\end{table*}

\subsection{Additional Span-Level Labeling Ablation}
To further validate the span-level semantic assignment component, we construct a dedicated polysemy evaluation set from HotpotQA. We collect text contexts containing spans with multiple semantics, and use LLM assistance to generate candidate meanings and gold semantic labels for these span occurrences. This dataset is designed to directly test whether the proposed semantic assignment algorithm can assign the correct meaning to each span in context, rather than only improving end-to-end retrieval scores.

We evaluate the proposed span-level semantic assignment algorithm on this dataset. The algorithm propagates sparse anchor labels over an occurrence embedding graph and uses a cross-encoder as semantic evidence for candidate scoring, unreliable-propagation veto, and fallback assignment. We compare this propagation-based assignment strategy against a cross-encoder baseline that independently labels every occurrence by selecting the highest-scoring candidate sense. Table~\ref{tab:span_ablation} summarizes the results of this diagnostic evaluation.

\begin{table}[t]
\centering
\small
\setlength{\tabcolsep}{6pt}
\caption{Span-level polysemy labeling results on HotpotQA}
\label{tab:span_ablation}
\begin{tabular}{lcc}
\toprule
Metric & Macro mean & Std. \\
\midrule
Algorithm overall accuracy & 93.7\% & 4.2\% \\
Propagation accuracy & 89.3\% & 17.6\% \\
Cross-encoder accuracy & 66.5\% & 25.7\% \\
\bottomrule
\end{tabular}
\end{table}

The results show that the proposed algorithm provides a substantial improvement over labeling each occurrence independently with the cross-encoder. The higher algorithm overall accuracy indicates that combining sparse anchor evidence, graph propagation, and cross-encoder verification is more reliable than using isolated candidate scoring alone. The propagation accuracy further shows that the occurrence embedding graph can correctly transfer semantic labels to nearby contexts, rather than merely serving as a wrapper around cross-encoder fallback.

The relatively large standard deviation of propagation accuracy reflects the difficulty of span-level polysemy across lexical anchors. Some anchors have compact and well-separated contextual regions, while others contain rare senses or semantically gradual transitions that are harder to propagate reliably. This behavior supports the design choice of combining graph propagation with cross-encoder verification and fallback: propagation handles locally coherent cases efficiently, while the verifier prevents uncertain or rare boundary cases from being forced into an unreliable propagated label.


\section{Conclusion}
In this paper, we proposed LiteSemRAG, a lightweight and LLM-frugal semantic-aware graph retrieval framework for Retrieval-Augmented Generation. LiteSemRAG constructs a heterogeneous semantic graph from contextual span embeddings, explicitly separating surface lexical forms from context-dependent semantic meanings to model token polysemy without LLM-centered graph construction or reasoning. The index preserves phrase-level evidence as part of graph construction, which helps maintain lexical precision while reducing unnecessary fragmentation of multiword expressions. The framework introduces dynamic semantic node construction based on anchor-guided semantic assignment, combining farthest-first anchor sampling, mutual-kNN label propagation, cross-encoder margin veto, and cross-encoder fallback, together with semantic co-occurrence retrieval that combines direct semantic evidence with bounded pairwise chunk support. For anchors requiring polysemy resolution, the assigned candidate-meaning labels provide an interpretable basis for inspecting and validating discovered senses. Experiments on three benchmark datasets show that LiteSemRAG achieves competitive or superior retrieval performance compared to LLM-based graph RAG systems while substantially improving indexing and querying efficiency through targeted indexing-time LLM use and encoder-driven retrieval.

%
% ---- Bibliography ----
%
% BibTeX users should specify bibliography style 'IEEEtran'.
% References will then be sorted and formatted in the correct style.
%
\bibliographystyle{IEEEtran}
\bibliography{references}

\end{document}
